\chapter{Manual de instalare și utilizare}
\pagestyle{fancy}

\section{Scopul manualului}

Acest capitol prezintă pașii necesari pentru instalarea, pornirea și utilizarea aplicației dezvoltate. Manualul este orientat către rularea locală a sistemului, în scop demonstrativ și experimental. Prin urmare, sunt descrise resursele necesare, comenzile de pornire, inițializarea datelor demonstrative și principalele interacțiuni disponibile în interfața web.

\section{Cerințe software și hardware}

Pentru rularea aplicației este necesar un sistem capabil să execute containere Docker și procese Node.js. Cerințele recomandate sunt prezentate în Tabelul \ref{tab:cerinte-instalare}.

\begin{table}[htbp]
\centering
\small
\begin{tabular}{|p{0.32\textwidth}|p{0.54\textwidth}|}
\hline
\textbf{Resursă} & \textbf{Cerință recomandată} \\
\hline
Sistem de operare & Windows 10/11 sau o distribuție Linux compatibilă cu Docker. \\
\hline
Docker & Docker Desktop sau Docker Engine, utilizat pentru pornirea serviciilor locale. \\
\hline
Node.js și npm & Versiune LTS recentă, necesară pentru aplicația Next.js și scripturile auxiliare. \\
\hline
Git & Necesar pentru obținerea proiectului din repository. \\
\hline
Browser web & Browser modern, de exemplu Google Chrome, Microsoft Edge sau Firefox. \\
\hline
Memorie RAM & Minimum 8 GB RAM recomandați pentru rularea simultană a serviciilor. \\
\hline
Spațiu pe disc & Spațiu liber pentru imaginile Docker, volumele locale și dependențele proiectului. \\
\hline
Conexiune internet & Necesară la prima instalare, pentru descărcarea imaginilor Docker și a pachetelor npm. \\
\hline
\end{tabular}
\caption{Cerințe software și hardware pentru rularea aplicației}
\label{tab:cerinte-instalare}
\end{table}

Python este necesar doar dacă serviciul de predicție ARIMA este rulat manual în afara containerelor Docker. În configurația locală recomandată, acesta este pornit prin Docker Compose.

\section{Configurarea mediului local}

După obținerea proiectului, terminalul trebuie deschis în directorul rădăcină al aplicației. Fișierul \texttt{docker-compose.yml} definește serviciile locale, iar fișierul \texttt{web-app/.env} conține variabilele utilizate de aplicația web și de worker-ul MQTT.

Pentru mediul demonstrativ local se pot utiliza valorile implicite din Tabelul \ref{tab:variabile-mediu}. Aceste valori sunt adecvate pentru rulare locală, dar nu trebuie folosite nemodificate într-un mediu de producție.

\begin{table}[htbp]
\centering
\small
\begin{tabular}{|p{0.34\textwidth}|p{0.50\textwidth}|}
\hline
\textbf{Variabilă} & \textbf{Rol} \\
\hline
\texttt{DATABASE\_URL} & Conexiunea aplicației la baza de date PostgreSQL pentru metadate. \\
\hline
\texttt{CHIRPSTACK\_DATABASE\_URL} & Conexiunea utilizată de scripturile de provisioning pentru baza ChirpStack. \\
\hline
\texttt{INFLUX\_URL} & Adresa locală a serviciului InfluxDB. \\
\hline
\texttt{INFLUX\_TOKEN} & Tokenul folosit pentru accesarea bucket-ului InfluxDB. \\
\hline
\texttt{INFLUX\_ORG} & Organizația InfluxDB configurată local. \\
\hline
\texttt{INFLUX\_BUCKET} & Bucket-ul în care sunt salvate citirile contoarelor. \\
\hline
\texttt{MQTT\_BROKER\_URL} & Adresa brokerului MQTT local. \\
\hline
\texttt{FORECAST\_SERVICE\_URL} & Adresa serviciului ARIMA, expus local pe portul \texttt{8001}. \\
\hline
\texttt{SIMULATOR\_API\_URL} & Adresa API-ului simulatorului LoRaWAN. \\
\hline
\texttt{JWT\_SECRET} & Cheia locală folosită pentru semnarea sesiunilor utilizatorilor. \\
\hline
\end{tabular}
\caption{Variabile de mediu importante pentru rularea locală}
\label{tab:variabile-mediu}
\end{table}

\section{Pornirea aplicației}

Pornirea mediului local se realizează în mai multe etape. Mai întâi se pornesc serviciile definite în Docker Compose și se pregătește aplicația web din directorul \texttt{web-app}.

\begin{verbatim}
docker compose up -d --build
cd web-app
npm install
npx prisma migrate deploy
npm run prisma:generate
\end{verbatim}

După inițializarea datelor demonstrative, descrisă în secțiunea următoare, aplicația web și worker-ul MQTT se pornesc prin:

\begin{verbatim}
cd web-app
npm run dev:all
\end{verbatim}

Comanda \texttt{npm run dev:all} pornește simultan aplicația web și procesul care recepționează mesajele MQTT. Pentru ca datele live să fie actualizate periodic în simulator, se deschide un al doilea terminal și se rulează:

\begin{verbatim}
cd web-app
npm run simulator:live
\end{verbatim}

Acest script nu înlocuiește simulatorul, ci actualizează payload-urile dispozitivelor prin mecanismul existent al simulatorului. După pornirea aplicației, se accesează interfața LWN Simulator la adresa \texttt{http://localhost:8000} și se pornește simularea. Aplicația web poate fi accesată în browser la adresa \texttt{http://localhost:3000}.

Serviciile locale principale sunt prezentate în Tabelul \ref{tab:servicii-porturi}.

\begin{table}[htbp]
\centering
\small
\begin{tabular}{|p{0.32\textwidth}|p{0.44\textwidth}|p{0.12\textwidth}|}
\hline
\textbf{Serviciu} & \textbf{Adresă locală} & \textbf{Rol} \\
\hline
Aplicația web & \texttt{http://localhost:3000} & Dashboard \\
\hline
ChirpStack & \texttt{http://localhost:8080} & Server LoRaWAN \\
\hline
InfluxDB & \texttt{http://localhost:8086} & Date istorice \\
\hline
LWN Simulator & \texttt{http://localhost:8000} & Simulare \\
\hline
Serviciu ARIMA & \texttt{http://localhost:8001} & Predicție \\
\hline
\end{tabular}
\caption{Servicii locale și porturi utilizate}
\label{tab:servicii-porturi}
\end{table}

\section{Inițializarea datelor demonstrative}

Pentru ca dashboard-ul să conțină utilizatori, contoare și date istorice, trebuie rulate scripturile de inițializare din directorul \texttt{web-app}:

\begin{verbatim}
cd web-app
npm run provision:demo-devices
npm run seed:history
\end{verbatim}

Scriptul \texttt{provision:demo-devices} creează sau actualizează conturile demonstrative și dispozitivele asociate acestora. Dispozitivele sunt create în simulator, în ChirpStack și în baza de date a aplicației, păstrând același identificator \texttt{devEui}. În plus, scriptul configurează profilurile ChirpStack necesare pentru tipurile de contoare utilizate în demonstrație, astfel încât payload-urile pentru electricitate, gaz, apă, încălzire și răcire să fie decodate corect. Scriptul \texttt{seed:history} generează citiri istorice pentru graficele de consum și pentru modulul ARIMA.

Pentru funcționarea fluxului live sunt necesare două acțiuni suplimentare: pornirea simulării din interfața LWN Simulator și rularea comenzii \texttt{npm run simulator:live} într-un terminal separat. Comanda trebuie lăsată activă pe durata demonstrației, deoarece actualizează periodic valorile transmise de dispozitive.

Conturile demonstrative sunt prezentate în Tabelul \ref{tab:conturi-demo}. 

\begin{table}[htbp]
\centering
\small
\begin{tabular}{|p{0.34\textwidth}|p{0.20\textwidth}|p{0.30\textwidth}|}
\hline
\textbf{Email} & \textbf{Parolă} & \textbf{Cod de revendicare} \\
\hline
\texttt{company.demo@example.com} & \texttt{Demo12345!} & \texttt{COMPANY-DEMO-2026} \\
\hline
\texttt{user1.demo@example.com} & \texttt{Demo12345!} & \texttt{USER1-DEMO-2026} \\
\hline
\texttt{user2.demo@example.com} & \texttt{Demo12345!} & \texttt{USER2-DEMO-2026} \\
\hline
\texttt{user3.demo@example.com} & \texttt{Demo12345!} & \texttt{USER3-DEMO-2026} \\
\hline
\end{tabular}
\caption{Conturi demonstrative disponibile în aplicație}
\label{tab:conturi-demo}
\end{table}

\section{Utilizarea aplicației}

\subsection{Autentificarea}

Utilizatorul accesează aplicația la adresa \texttt{http://localhost:3000} și introduce adresa de email și parola contului. Pentru demonstrarea scenariului de companie se poate utiliza contul \texttt{company.demo@example.com}.

\subsection{Prezentarea generală a flotei}

Pagina de prezentare generală oferă o imagine agregată asupra dispozitivelor asociate contului. Utilizatorul poate observa numărul total de contoare, categoriile de utilități, costurile estimate și starea fluxului live. 

\begin{figure}[h!]
\centering
\includegraphics[width=0.75\textwidth]{imagini/02_overview.png}
\caption{Pagina de prezentare generală a flotei}
\label{fig:manual-overview}
\end{figure}

\subsection{Administrarea dispozitivelor}

Pagina de dispozitive permite vizualizarea inventarului de contoare. Utilizatorul poate căuta un dispozitiv după nume sau \texttt{devEui}, poate filtra lista după stare și poate accesa acțiuni precum vizualizarea detaliilor sau editarea metadatelor. Pentru flote mai mari, lista este paginată, astfel încât interfața să rămână ușor de utilizat.

\begin{figure}[h!]
\centering
\includegraphics[width=0.75\textwidth]{imagini/03_devices.png}
\caption{Pagina de administrare a dispozitivelor}
\label{fig:manual-devices}
\end{figure}

În cazul unui cont nou, utilizatorul poate folosi un cod de revendicare pentru a asocia dispozitivele pregătite anterior. Codul se introduce în panoul dedicat din pagina de dispozitive.

\subsection{Vizualizarea unui contor individual}

Pagina contorului individual prezintă informațiile detaliate pentru dispozitivul selectat. Utilizatorul poate consulta citirile recente, profilul de consum, costul estimat și predicția ARIMA pentru intervalul următor. Selectorul de dispozitive permite schimbarea contorului analizat fără revenirea la pagina de inventar.

\begin{figure}[h!]
\centering
\includegraphics[width=0.75\textwidth]{imagini/04_meter.png}
\caption{Pagina contorului individual și prognoza ARIMA}
\label{fig:manual-meter}
\end{figure}

\subsection{Analiza costurilor}

Pagina de facturare prezintă costurile estimate pentru dispozitive și pentru categoriile de utilități. Valorile afișate au rol orientativ și sunt calculate pe baza tarifelor definite în aplicație și a consumurilor disponibile. Această pagină este utilă pentru identificarea contoarelor sau utilităților care contribuie cel mai mult la costul total.

\begin{figure}[h!]
\centering
\includegraphics[width=0.75\textwidth]{imagini/05_billing.png}
\caption{Pagina de analiză a costurilor}
\label{fig:manual-billing}
\end{figure}

\subsection{Vizualizarea pe hartă}

Modul de hartă este disponibil din pagina de prezentare generală. Acesta afișează contoarele pe baza coordonatelor geografice configurate. Utilizatorul poate filtra dispozitivele după tipul utilității sau stare și poate selecta un contor pentru a vedea informațiile sale principale.

\begin{figure}[h!]
\centering
\includegraphics[width=0.75\textwidth]{imagini/07_map.png}
\caption{Vizualizarea geografică a dispozitivelor}
\label{fig:manual-map}
\end{figure}
